\section{Related Works}

This section surveys prior work most closely related to our setting.

\subsection{Financial Network Risk Modelling}
In traditional finance, credit risk modelling typically spans both entity-level default modelling and system-level network perspectives. Classical structural and reduced-form approaches focus on estimating default likelihoods for individual institutions, whereas network contagion models study how distress can cascade through interconnected balance sheets. For instance, Dolfin et al. use network epidemic models to study credit contagion, showing that dense interconnections and overlapping portfolios can magnify systemic credit risk \cite{math7080713}.

In decentralised finance, credit risk takes a different form because lending is commonly over-collateralised and enforced by smart contracts. Bertomeu et al. introduce parsimonious aggregate risk measures for DeFi lending protocols based only on total deposits and borrowings, and find that systemic fragility rose sharply around mid-2021 during periods of crypto-market stress \cite{BERTOMEU2024105321}. Bridging conceptual links between TradFi and DeFi, Aufiero et al. review mechanisms of systemic risk and argue that, while familiar ingredients (leverage, liquidity shocks, correlated exposures) exist in both systems, DeFi’s algorithmic execution and composability introduce distinctive channels through which shocks may propagate \cite{aufiero2025mappingmicroscopicsystemicrisks}.

\subsection{Machine Learning Approaches}

Machine learning methods are increasingly used for systemic-risk detection and contagion analysis \cite{BALMASEDA2023200240, gonon2025computingsystemicriskmeasures}. In particular, graph neural networks (GNNs) exploit relational structure to strengthen predictive performance on networked risk problems \cite{KipfWelling2017, veličković2018graphattentionnetworks}. Balmaseda et al. demonstrate that GNN-based models can substantially outperform conventional ML when classifying the systemic importance of banks in simulated financial networks \cite{BALMASEDA2023200240}. Gonon et al. combine GNNs with explicit interbank liability networks to compute systemic risk measures defined via the minimum capital required to stabilise the system \cite{gonon2025computingsystemicriskmeasures}. Their formulation effectively captures contagion channels (e.g., Eisenberg--Noe-style clearing) by embedding graph-structured information directly into the risk aggregation procedure.

\subsection{Foundation Models}

Foundation models are increasingly influencing finance and economics by shifting from narrowly tuned, task-specific models towards reusable representations learned from broad data. Rather than training a bespoke model per downstream objective, recent work pretrains large neural networks---predominantly transformer-based architectures \cite{VaswaniEtAl2017Attention} and related variants---on massive corpora, then transfers these models across tasks via prompting or lightweight adaptation \cite{AnsariEtAl2024Chronos,RasulEtAl2023LagLlama,DasEtAl2024TimesFM}. Alongside this, an emerging policy-oriented literature uses large language models to extract quantitative signals from narrative sources (e.g., reports, releases, and news) and to integrate them into nowcasting pipelines, suggesting that relatively small yet information-rich text inputs can improve real-time forecasting when paired with standard indicators \cite{deBondtSun2025ChatGPTNowcast,KwonEtAl2024BISLLMprimer,CarrieroEtAl2024MacroLLM}. 

Foundation-model ideas are also being extended to graph-structured inputs by augmenting self-attention with graph-aware structural encodings, illustrating the viability of large pretrained models for networks that resemble webs of financial contracts or cross-asset relationships; Graphormer provides a representative example \cite{Ying2021Graphormer}. For tabular settings, the Graph Prior-data Fitted Network (GraphPFN) \cite{eremeev2025turningtabularfoundationmodels} is pretrained on millions of synthetic tabular tasks and reports strong out-of-the-box performance across a range of benchmarks.